\documentclass[12pt,a4paper]{article}
\usepackage{cmap} % Makes the PDF copiable. See http://tex.stackexchange.com/a/64198/25761
\usepackage[T1]{fontenc}
\usepackage[brazil]{babel}
\usepackage[utf8]{inputenc}
\usepackage{amsmath}
\usepackage{amsfonts}
\usepackage{amssymb}
\usepackage{amsthm}
\usepackage{textcomp} % \degree
\usepackage{gensymb} % \degree
\usepackage[usenames,svgnames,dvipsnames]{xcolor}
\usepackage{hyperref}
\usepackage{graphicx}
\usepackage[margin=2cm]{geometry}

\hypersetup{
    colorlinks = true,
    allcolors = {blue}
}

\newcommand*\sen[1]{\operatorname{sen}{#1}}
\newcommand*\abs[1]{\left|#1\right|}

\newcommand*\tipo{PROVA I}
\newcommand*\turma{MAT102-02U}
\newcommand*\disciplina{CDI0001}
\newcommand*\eu{Helder G. G. de Lima}
\newcommand*\data{16/09/2015}

\author{\eu}
\title{\tipo - \disciplina}
\date{\data}

\begin{document}
\thispagestyle{empty}
\newgeometry{margin=2cm,bottom=0.5cm}
\begin{center}
\includegraphics[width=9.0cm]{marca}
\noindent\begin{tabular}{l c c r}
  \textbf{\disciplina}
& \textbf{\tipo}
& \textbf{\data}
& \textbf{\turma}
\end{tabular}
\\ Prof. \eu\footnote{
Este é um material de acesso livre distribuído sob os termos da licença \href{https://creativecommons.org/licenses/by-sa/4.0/deed.pt_BR}{Creative Commons BY-SA 4.0}.}
\end{center}

\noindent Nome do(a) aluno(a): \rule{13cm}{0.01cm}

%\section*{Instruções}
\begin{center}\fbox{
\begin{minipage}{14cm}

{\footnotesize
\begin{itemize}
\renewcommand{\theenumi}{\Roman{enumi}}
\item Identifique-se em todas as folhas.
\item Mantenha o celular e os demais equipamentos eletrônicos desligados durante a prova.
\item Justifique cada resposta com cálculos ou argumentos baseados na teoria estudada.
\item Escolha \textsc{\textbf{uma}} das 6 questões para \textsc{\textbf{não}} fazer (ela não será corrigida): \framebox(30,10){}
\end{itemize}
}

\end{minipage}
}
\end{center}


%\section*{Questões}

\begin{enumerate}
\item Seja $g: D \to \mathbb{R}$ a função dada por $g(x) = \ln{|x + 3|} + \ln{|x - 3|} + x^2\sqrt{x^2-1}$. Explique:
\begin{enumerate}
\item (0,5) Qual é o domínio $D$ em que $g$ está definida?
\item (0,5) A função $g$ é par, ímpar ou nenhuma das duas coisas?
\item (0,5) A função $g$ é contínua nos pontos em que está definida? Por quê?
\item (0,5) Quanto é $\lim\limits_{x\to3^+} g(x)$?
\end{enumerate}

\item Seja $f(x) = \ln(x - 2)$ e considere $\varepsilon > 0$.
\begin{enumerate}
\item (1,5) Obtenha $\delta > 0$ (em função de $\varepsilon$) tal que ``se $0 < |x - 3| < \delta$ então $|f(x) - 0| < \varepsilon$''.
\item (0,5) O que isso significa em termos de limites?
\end{enumerate}

\item (2,0) Determine um intervalo $[n, n+1]$ contendo (pelo menos) uma solução positiva da equação $x^5 - 10x = 300$. O que garante que há realmente uma solução no intervalo encontrado?

\item (2,0) Determine se para algum valor de $k$ a função a seguir é contínua em $x = -2$:
\[f(x)=\left\{\begin{array}{ll}
\dfrac{\sen(x+2)}{x^2-x-6}, & \text{se}\quad x<-2\\
k, & \text{se}\quad x = -2 \\
\dfrac{1}{-5 + e^{-1/(x + 2)}}, & \text{se}\quad x > -2.
\end{array}\right.
\]

\item Determine, se existirem:
\begin{enumerate}
\item (1,0) $\lim\limits_{x\to a} \dfrac{\sen(x)-\sen(a)}{x-a}$
\item (1,0) $\lim\limits_{x\to +\infty} \dfrac{\cos(x^2)}{e^x}$
\end{enumerate}

\item Calcule, se existirem:
\begin{enumerate}
\item (1,0) $\lim\limits_{x\to 0} \dfrac{e^{6x}-1}{1-e^{3x}} $
\item (1,0) $\lim\limits_{x\to +\infty} x(\sqrt{x^2 + 10} - x)$
\end{enumerate}
\end{enumerate}

\newpage
\restoregeometry
\section*{Respostas e observações}
\begin{enumerate}

\item \textit{ Seja $g: D \to \mathbb{R}$ a função dada por $g(x) = \ln{|x + 3|} + \ln{|x - 3|} + x^2\sqrt{x^2-1}$. Explique:}
\begin{enumerate}
\item \textit{ Qual é o domínio $D$ em que $g$ está definida?}

\begin{align*}
D
& = \operatorname{Dom}(g) \\
& = \left\{ x \in \mathbb{R} | \abs{x + 3} > 0 \text{ e } \abs{x-3} > 0 \text{ e } x^2-1 \geq 0 \right\} \\
& = \left\{ x \in \mathbb{R} | \abs{x + 3} \neq 0 \text{ e } \abs{x-3} \neq 0 \text{ e } x^2 \geq 1 \right\} \\
& = \left\{ x \in \mathbb{R} | x+3 \neq 0 \text{ e } x-3 \neq 0 \text{ e } \abs{x} \geq 1 \right\} \\
& = \left\{ x \in \mathbb{R} | x \neq -3 \text{ e } x\neq 3 \text{ e } \abs{x} \geq 1 \right\} \\
& = \mathbb{R} \setminus \left( \{-3, 3 \} \cup (-1, 1) \right)\\
& = (-\infty, -3) \cup ( -3, -1] \cup [ 1, 3) \cup (1, +\infty)
\end{align*}

\item \textit{ A função $g$ é par, ímpar ou nenhuma das duas coisas?}

Tem-se
\begin{align*}
g(-x)
& = \ln{|(-x) + 3|} + \ln{|(-x) - 3|} + (-x)^2\sqrt{(-x)^2-1} \\
& = \ln{|-(x - 3)|} + \ln{|-(x + 3)|} + x^2\sqrt{x^2-1} \\
& = \ln{|x - 3|} + \ln{|x + 3|} + x^2\sqrt{x^2-1} \\
& = g(x).
\end{align*}
Isso significa que $g$ é uma função par.

\item \textit{ A função $g$ é contínua nos pontos em que está definida? Por quê?}

Uma vez que raízes quadradas, polinômios e logaritmos são funções contínuas, e que a soma, subtração, multiplicação e composição de funções contínuas resulta em funções contínuas, pode-se concluir que a função $g$ também é contínua (pois é uma soma/composição de funções contínuas).

\item \textit{ Quanto é $\lim\limits_{x\to3^+} g(x)$?}
Levando em conta que
\[
\lim\limits_{x\to3^+} \ln{|x + 3|} + x^2\sqrt{x^2-1}
= \ln{|3 + 3|} + 3^2\sqrt{3^2-1}
= \ln{6} + 9\sqrt{8} \in \mathbb{R}
\]
e que
\[
\lim\limits_{x\to3^+} \ln{\abs{x - 3}}
= \lim\limits_{u\to 0^+} \ln{\abs{u}}
= \lim\limits_{u\to 0^+} \ln{u}
= -\infty,
\]
conclui-se que
\[
\lim\limits_{x\to 3^+} g(x) = - \infty.
\]

\end{enumerate}

\item \textit{ Seja $f(x) = \ln(x - 2)$ e considere $\varepsilon > 0$.}
\begin{enumerate}
\item \textit{ Obtenha $\delta > 0$ (em função de $\varepsilon$) tal que ``se $0 < |x - 3| < \delta$ então $|f(x) - 0| < \varepsilon$''.}

Observe que
\begin{align*}
\abs{f(x) - 0} < \varepsilon
& \Leftrightarrow \abs{ \ln(x - 2) } < \varepsilon
  \Leftrightarrow -\varepsilon < \ln(x - 2) < \varepsilon
  \Leftrightarrow e^{-\varepsilon} < e^{\ln(x - 2)} < e^\varepsilon \\
& \Leftrightarrow e^{-\varepsilon} < x - 2 < e^\varepsilon
  \Leftrightarrow e^{-\varepsilon} -1 < x - 3 < e^\varepsilon -1.
\end{align*}
Portanto, para um $\varepsilon>0$ dado, basta escolher qualquer $\delta \leq \min\{ \abs{e^{-\varepsilon} -1 }, e^\varepsilon -1\}$ e será verdade que ``se $0 < \abs{x - 3} < \delta$ então $\abs{f(x) - 0} < \varepsilon$''.



\item \textit{ O que isso significa em termos de limites?}

O significado da propriedade anterior é que a função $f(x)$ se aproxima de zero quando $x$ tende a $3$, isto é, $\lim\limits_{x\to3} f(x) = 0$.

\end{enumerate}

\item \textit{ Determine um intervalo $[n, n+1]$ contendo (pelo menos) uma solução positiva da equação $x^5 - 10x = 300$. O que garante que há realmente uma solução no intervalo encontrado?}

Seja $f(x) = x^5 - 10x$. Atribuindo alguns valores à variável $x$ (positivos, já que o enunciado não pede soluções negativas), observa-se que
\begin{align*}
f(1) & = 1^5 - 10 \cdot 1 = 1 - 10 = -9 < 300 \\
f(2) & = 2^5 - 10 \cdot 2 = 32 - 20 = 12 < 300 \\
\bf{ f(3) } & = 3^5 - 10 \cdot 3 = 243 - 30 = 213 \bf{ < 300 }\\
\bf{ f(4) } & = 4^5 - 10 \cdot 4 = 1024 - 40 = 984 \bf{ > 300 }
\end{align*}
Como $f$ é uma função contínua (é um polinômio) e $f(3) < 300 < f(4)$, o teorema do valor intermediário garante que existe algum $c \in (3,4)$ tal que $f(c) = c^5 - 10 \cdot c = 300$, isto é, $c$ é uma solução da equação indicada.

\item \textit{ Determine se para algum valor de $k$ a função a seguir é contínua em $x = -2$:}
\[f(x)=\left\{\begin{array}{ll}
\dfrac{\sen(x+2)}{x^2-x-6}, & \text{se}\quad x<-2\\
k, & \text{se}\quad x = -2 \\
\dfrac{1}{-5 + e^{-1/(x + 2)}}, & \text{se}\quad x > -2.
\end{array}\right.
\]
Para $x<-2$ temos o seguinte limite:
\begin{align*}
\lim\limits_{x\to -2^-} f(x)
& = \lim\limits_{x\to -2^-} \frac{\sen(x+2)}{x^2-x-6}
  = \lim\limits_{x\to -2^-} \frac{\sen(x+2)}{(x+2)(x-3)}
  = \lim\limits_{u\to 0^-} \frac{\sen(u)}{u}\frac{1}{(u-2)-3}
  = \frac{-1}{5}.
\end{align*}

Para $x>-2$, temos
\begin{align*}
\lim\limits_{x\to -2^+} f(x)
& = \lim\limits_{x\to -2^+} \frac{1}{-5 + e^{-1/(x + 2)}}
  = \lim\limits_{u\to 0^+} \frac{1}{-5 + e^{-1/u}}
  = \lim\limits_{v\to -\infty} \frac{1}{-5 + e^v}
  = \frac{1}{-5 + 0}
  = \frac{-1}{5}.
\end{align*}

Consequentemente, $f$ é contínua em $x=-2$ se, e somente se, $k = -1/5$, já que neste caso $\lim\limits_{x\to -2} f(x) = -1/5 = k = f(-2)$.

\item \textit{ Determine, se existirem:}
\begin{enumerate}
\item
\begin{align*}
\lim\limits_{x\to a} \frac{\sen x-\sen a}{x-a}
& = \lim\limits_{u \to 0} \frac{\sen(a+u)-\sen a}{u}\\
& = \lim\limits_{u \to 0} \frac{\sen a \cos u +\sen u\cos a-\sen a}{u}\\
& = \lim\limits_{u \to 0} \frac{(\cos u - 1)}{u}\sen a + \frac{\sen u}{u}\cos a\\
& = 0 \cdot \sen a + 1 \cdot \cos a
  = \cos a.\\
\end{align*}

\item $\lim\limits_{x\to +\infty} \dfrac{\cos(x^2)}{e^x} = \lim\limits_{x\to +\infty} \dfrac{1}{e^x} \cdot \cos(x^2) = 0$, pois $\lim\limits_{x\to +\infty} \dfrac{1}{e^x} = 0$ e $-1 \leq \cos(x^2) \leq 1$, isto é, a função limitada $\cos(x^2)$ está sendo multiplicada por uma que tende a zero, o que implica que o produto também tende a zero.

\end{enumerate}

\item \textit{ Calcule, se existirem:}
\begin{enumerate}
\item $\lim\limits_{x\to 0} \dfrac{e^{6x}-1}{1-e^{3x}}$

\begin{align*}
\lim\limits_{x\to 0} \frac{e^{6x}-1}{1-e^{3x}}
& = \lim\limits_{x\to 0} \frac{6x \cdot \frac{e^{6x}-1}{6x}}{3x \cdot \frac{1-e^{3x}}{3x}}
  = 2\lim\limits_{x\to 0} \frac{\frac{e^{6x}-1}{6x}}{(-1)\frac{e^{3x}-1}{3x}}
  = (-2)\frac{\lim\limits_{x\to 0} \frac{e^{6x}-1}{6x}}{\lim\limits_{x\to 0} \frac{e^{3x}-1}{3x}} \\
& = (-2)\frac{\lim\limits_{u\to 0} \frac{e^{u}-1}{u}}{\lim\limits_{v\to 0} \frac{e^{v}-1}{v}}
  = (-2) \cdot \frac{1}{1}
  = -2
\end{align*}

\item $\lim\limits_{x\to +\infty} x(\sqrt{x^2 + 10} - x)$

\begin{align*}
\lim\limits_{x\to +\infty} x(\sqrt{x^2 + 10} - x)
& = \lim\limits_{x\to +\infty} x(\sqrt{x^2 + 10} - x) \cdot \frac{\sqrt{x^2 + 10} + x}{\sqrt{x^2 + 10} + x} \\
& = \lim\limits_{x\to +\infty} \frac{x((x^2 + 10) - x^2)}{\sqrt{x^2 + 10} + x}
  = \lim\limits_{x\to +\infty} \frac{10 x}{|x|\sqrt{1 + 10/x^2} + x} \\
& = \lim\limits_{x\to +\infty} \frac{10 x}{x (\sqrt{1 + 10/x^2} + 1)}
  = \frac{10}{\sqrt{1 + 0} + 1}
  = 5
\end{align*}


\end{enumerate}
\end{enumerate}

\end{document}
